\subsection{Frontend-Datenoperationen}

Die Patientenansicht wurde um schreibende Funktionen erweitert. Neben der bisherigen Anzeige von Patienten, Diagnosen und Buchungen können nun neue Patienten angelegt sowie fachliche Aktionen zu bestehenden Patienten ausgeführt werden (siehe Abbildungen~\ref{fig:frontend-add-patient},~\ref{fig:frontend-patient-actions-1} und~\ref{fig:frontend-patient-actions-2}). Für diese Operationen wurde eine einheitliche Verarbeitung von JSON-Anfragen ergänzt, sodass erfolgreiche Änderungen und Validierungsfehler konsistent in der Oberfläche abgebildet werden.

\subsubsection{Patienten- und Buchungsaktionen}

Für die Patientenerstellung wurde ein Dialog ergänzt, der die Stammdaten der Person strukturiert erfasst (siehe Abbildung~\ref{fig:frontend-add-patient}). Meldet das Backend mögliche Übereinstimmungen mit bestehenden Personen zurück, werden diese im Frontend angezeigt und können für die Patientenerstellung weiterverwendet werden. Dadurch wird die Trennung zwischen Person und Patient in der Oberfläche sichtbar: Eine Person muss nicht doppelt gespeichert werden, sondern kann bei Bedarf eine zusätzliche Patientenrolle erhalten.

In der Patientendetailansicht wurden Aktionen für Zimmerbuchungen ergänzt. Eine Buchung kann neu angelegt, ein Patient in ein anderes Zimmer verlegt oder entlassen werden (siehe Abbildungen~\ref{fig:frontend-patient-actions-1} und~\ref{fig:frontend-patient-actions-2}). Die Verlegung und Entlassung werden nur angeboten, wenn eine aktuelle Buchung vorhanden ist. Damit orientiert sich die Oberfläche an der Datenstruktur der Buchungen, in der der Aufenthalt eines Patienten über Zeitraum, Zimmer und Status modelliert wird. Nach einer erfolgreichen Aktion werden die Patientendetails erneut geladen, damit die Anzeige dem gespeicherten Zustand entspricht.

\subsubsection{Statusfilter für Patienten}

Die Patientenliste wurde um einen Mehrfachfilter für den Buchungsstatus erweitert (siehe Abbildung~\ref{fig:frontend-filter}). Ausgewählt werden kann, ob Patienten aktuell eingecheckt sind oder eine kommende Buchung besitzen. Technisch werden mehrere Statuswerte als wiederholbare Query-Parameter übertragen.

Die Auswertung erfolgt nicht über Personenattribute, sondern über die Beziehung zwischen Patient und Buchung. Für aktuell eingecheckte Patienten wird eine gültige Buchung mit dem Status \texttt{CHECKED\_IN} gesucht. Für kommende Aufenthalte werden zukünftige Buchungen mit den Statuswerten \texttt{PENDING} oder \texttt{CONFIRMED} berücksichtigt. Dadurch lässt sich die relationale Buchungsinformation mit den bestehenden Patientenfiltern kombinieren, ohne die Einschränkungen der verschlüsselten Personenfelder zu umgehen.

\subsubsection{Diagnosen, Medikamente und Dosierungen}

Für Patienten können nun Diagnosen direkt aus der Detailansicht heraus erfasst werden (siehe Abbildung~\ref{fig:frontend-add-diagnosis}). Der Dialog lädt Ärzte und Medikamente als Auswahlwerte und erfasst zusätzlich die Dosierung. Die Speicherung folgt der fachlichen Abhängigkeit der Daten: Zuerst wird die Dosierung erzeugt, danach die Medikation mit Bezug auf Medikament und Dosierung, anschließend die Diagnose mit Bezug auf Patient, Arzt und Medikation.

Diese Reihenfolge verhindert, dass abhängige Datensätze ohne ihre Referenzen entstehen. Die Werte für Dosierungsmenge und Frequenz werden numerisch erfasst und backendseitig vor der Speicherung geprüft. Zusätzlich können bestehende Diagnosen beendet werden. Dabei wird der Lebenszyklus der Diagnose geändert, ohne den historischen Datensatz zu entfernen.
